% Options for packages loaded elsewhere
\PassOptionsToPackage{unicode}{hyperref}
\PassOptionsToPackage{hyphens}{url}
\PassOptionsToPackage{dvipsnames,svgnames,x11names}{xcolor}
\documentclass[
  11pt,
]{article}
\usepackage{xcolor}
\usepackage[margin=25mm]{geometry}
\usepackage{amsmath,amssymb}
\setcounter{secnumdepth}{-\maxdimen} % remove section numbering
\usepackage{iftex}
\ifPDFTeX
  \usepackage[T1]{fontenc}
  \usepackage[utf8]{inputenc}
  \usepackage{textcomp} % provide euro and other symbols
\else % if luatex or xetex
  \usepackage{unicode-math} % this also loads fontspec
  \defaultfontfeatures{Scale=MatchLowercase}
  \defaultfontfeatures[\rmfamily]{Ligatures=TeX,Scale=1}
\fi
\usepackage{lmodern}
\ifPDFTeX\else
  % xetex/luatex font selection
\fi
% Use upquote if available, for straight quotes in verbatim environments
\IfFileExists{upquote.sty}{\usepackage{upquote}}{}
\IfFileExists{microtype.sty}{% use microtype if available
  \usepackage[]{microtype}
  \UseMicrotypeSet[protrusion]{basicmath} % disable protrusion for tt fonts
}{}
\makeatletter
\@ifundefined{KOMAClassName}{% if non-KOMA class
  \IfFileExists{parskip.sty}{%
    \usepackage{parskip}
  }{% else
    \setlength{\parindent}{0pt}
    \setlength{\parskip}{6pt plus 2pt minus 1pt}}
}{% if KOMA class
  \KOMAoptions{parskip=half}}
\makeatother
\usepackage{listings}
\newcommand{\passthrough}[1]{#1}
\lstset{defaultdialect=[5.3]Lua}
\lstset{defaultdialect=[x86masm]Assembler}
\usepackage{longtable,booktabs,array}
\usepackage{calc} % for calculating minipage widths
% Correct order of tables after \paragraph or \subparagraph
\usepackage{etoolbox}
\makeatletter
\patchcmd\longtable{\par}{\if@noskipsec\mbox{}\fi\par}{}{}
\makeatother
% Allow footnotes in longtable head/foot
\IfFileExists{footnotehyper.sty}{\usepackage{footnotehyper}}{\usepackage{footnote}}
\makesavenoteenv{longtable}
\setlength{\emergencystretch}{3em} % prevent overfull lines
\providecommand{\tightlist}{%
  \setlength{\itemsep}{0pt}\setlength{\parskip}{0pt}}
\emergencystretch=3em
\AtBeginDocument{\Urlmuskip=0mu plus 2mu}
\lstset{breaklines=true,basicstyle=\small\ttfamily,columns=fullflexible,keepspaces=true,showstringspaces=false}
\setlength{\tabcolsep}{4pt}
\usepackage{iftex}
\ifXeTeX
  \usepackage{newunicodechar}
  \makeatletter
  \newcommand{\DeclareUnicodeCharacter}[2]{%
    \begingroup
    \uccode`\~="#1\relax
    \uppercase{\def\temp{~}}%
    \expandafter\endgroup
    \expandafter\newunicodechar\expandafter{\temp}{#2}%
  }
  \makeatother
\fi
\ifdefined\DeclareUnicodeCharacter
\DeclareUnicodeCharacter{2192}{\ensuremath{\to}}
\DeclareUnicodeCharacter{2190}{\ensuremath{\leftarrow}}
\DeclareUnicodeCharacter{2194}{\ensuremath{\leftrightarrow}}
\DeclareUnicodeCharacter{21D2}{\ensuremath{\Rightarrow}}
\DeclareUnicodeCharacter{21A6}{\ensuremath{\mapsto}}
\DeclareUnicodeCharacter{2212}{\ensuremath{-}}
\DeclareUnicodeCharacter{2264}{\ensuremath{\leq}}
\DeclareUnicodeCharacter{2265}{\ensuremath{\geq}}
\DeclareUnicodeCharacter{2260}{\ensuremath{\neq}}
\DeclareUnicodeCharacter{2113}{\ensuremath{\ell}}
\DeclareUnicodeCharacter{03A3}{\ensuremath{\Sigma}}
\DeclareUnicodeCharacter{03A6}{\ensuremath{\Phi}}
\DeclareUnicodeCharacter{0393}{\ensuremath{\Gamma}}
\DeclareUnicodeCharacter{221E}{\ensuremath{\infty}}
\DeclareUnicodeCharacter{2208}{\ensuremath{\in}}
\DeclareUnicodeCharacter{2209}{\ensuremath{\notin}}
\DeclareUnicodeCharacter{211A}{\ensuremath{\mathbb{Q}}}
\DeclareUnicodeCharacter{211D}{\ensuremath{\mathbb{R}}}
\DeclareUnicodeCharacter{2124}{\ensuremath{\mathbb{Z}}}
\DeclareUnicodeCharacter{2115}{\ensuremath{\mathbb{N}}}
\DeclareUnicodeCharacter{2200}{\ensuremath{\forall}}
\DeclareUnicodeCharacter{2203}{\ensuremath{\exists}}
\DeclareUnicodeCharacter{2286}{\ensuremath{\subseteq}}
\DeclareUnicodeCharacter{00D7}{\ensuremath{\times}}
\DeclareUnicodeCharacter{03B4}{\ensuremath{\delta}}
\DeclareUnicodeCharacter{2202}{\ensuremath{\partial}}
\DeclareUnicodeCharacter{03B2}{\ensuremath{\beta}}
\DeclareUnicodeCharacter{03B1}{\ensuremath{\alpha}}
\DeclareUnicodeCharacter{03C6}{\ensuremath{\phi}}
\DeclareUnicodeCharacter{03B3}{\ensuremath{\gamma}}
\DeclareUnicodeCharacter{00B9}{\ensuremath{^1}}
\DeclareUnicodeCharacter{00B2}{\ensuremath{^2}}
\DeclareUnicodeCharacter{2070}{\ensuremath{^0}}
\DeclareUnicodeCharacter{2080}{\ensuremath{_0}}
\DeclareUnicodeCharacter{2081}{\ensuremath{_1}}
\DeclareUnicodeCharacter{2082}{\ensuremath{_2}}
\fi
\usepackage{bookmark}
\IfFileExists{xurl.sty}{\usepackage{xurl}}{} % add URL line breaks if available
\urlstyle{same}
\hypersetup{
  pdftitle={PFP VI: Retained Failure Evidence and Controlled Reuse},
  pdfauthor={Jeremy H. Carroll},
  colorlinks=true,
  linkcolor={Maroon},
  filecolor={Maroon},
  citecolor={Blue},
  urlcolor={Blue},
  pdfcreator={LaTeX via pandoc}}

\title{PFP VI: Retained Failure Evidence and Controlled Reuse}
\author{Jeremy H. Carroll}
\date{2026-10-02 · Version 0.1.0 · Systems methods preprint · First-publication candidate}

\AtBeginEnvironment{longtable}{\small}
\begin{document}
\makeatletter
\begin{center}
{\LARGE \@title\par}
\vspace{0.8em}
{\normalsize \@author\par}
\vspace{0.5em}
{\small \@date\par}
\end{center}
\makeatother

Code baseline:
\texttt{\nolinkurl{52cb5457f6c485ba548cfd292f6012a941ed22cf}}.

Publication package reviewed from
\texttt{\nolinkurl{2830d21f695e28dbfcf4f07e18a66da948b4a65b}}. The
implementation audit remains historical. The default branch's sink has
subsequently changed its float encoding and removed the RFC 8785 claim
from its encoder docstring. Those later changes are not imported into
the baseline argument or its test results. Version 0.1.0 denotes this
first public preprint candidate; earlier 1.x labels denote internal
drafts.

\hypertarget{abstract}{%
\subsection{Abstract}\label{abstract}}

A rejected proposal can explain a defect, reveal a mistaken assumption,
or supply a useful negative test. Keeping it also creates new
obligations: the record must remain interpretable, its integrity must be
checkable, and future reuse must not turn rejection into authority. This
paper examines those obligations in the Proposal Fidelity Protocol
(PFP), previously called the Symbol Game. The implementation
contribution is a source-audited JSONL failure ledger that records
structured refusals, chains entry digests, verifies restart state, and
emits a count-and-head receipt of its current in-memory state. We
describe exactly which modifications its checks detect and why an
independently retained anchor is needed when an attacker can rewrite the
whole file. A worked two-record example separates local replay, snapshot
comparison, and authorized retention. We then specify a controlled reuse
procedure using existing detector checks for benign examples,
trigger-based diagnostic downgrade thresholds, invalid patterns, and
affinity decay. A retained catalog of 61 harvested detector candidates
is described only as a data file; no code at the audited commit applies
its fixtures. The record store is logically append-only through its API;
it does not provide immutable media, authenticated provenance,
concurrent-writer safety, or a power-loss durability theorem. No
improvement in model accuracy, cost, or discovery yield is inferred from
retaining records.

\hypertarget{why-preserve-a-failed-attempt}{%
\subsection{1. Why preserve a failed
attempt?}\label{why-preserve-a-failed-attempt}}

Suppose a candidate code patch edits a protected file, or a proposed
mathematical witness refers to the wrong carrier. A host refuses the
transition. If the refusal is recorded only as ``failed,'' the next
investigator cannot tell which proposal was evaluated, which condition
failed, or whether the same failure has already been reproduced.

A useful failure record keeps enough authorized information to
reconstruct the attempt. That record can support a regression test and a
later comparison of detectors. It is still evidence about an event. It
does not establish that the proposal is false under every
interpretation, or that a recurring phrase is a reliable detector of
future errors.

This paper contributes a concrete storage contract, a retained-anchor
integrity argument, and a bounded procedure for promoting a stored
example into a proposed reusable rule. The procedure distinguishes three
actions: retaining an event, reproducing its failure, and admitting a
detector for a declared scope. None of these actions follows
automatically from the previous one.

The audited implementation is fixed at commit
\texttt{\nolinkurl{52cb5457f6c485ba548cfd292f6012a941ed22cf}}. The
main sources are
\href{https://github.com/dbccpage/the_symbol_game/blob/52cb5457f6c485ba548cfd292f6012a941ed22cf/symbol_game/infrastructure/persistence/cemetery_sink.py}{\passthrough{\lstinline!CemeterySink!}},
its
\href{https://github.com/dbccpage/the_symbol_game/blob/52cb5457f6c485ba548cfd292f6012a941ed22cf/symbol_game/domain/entities/cemetery_entry.py}{domain
record types}, and the
\href{https://github.com/dbccpage/the_symbol_game/blob/52cb5457f6c485ba548cfd292f6012a941ed22cf/symbol_game/tests/test_cemetery_sink.py}{storage
regression suite}. The inherited ``cemetery'' name means a store for
rejected attempts. The concrete data structure is an ordered predecessor
chain with a header, not a general graph of arbitrary dependencies. A
chain is a special kind of directed acyclic graph when edges are
oriented from each new entry to its earlier predecessor; that structural
observation supplies neither arbitrary dependency edges nor a semantic
mining algorithm.

\hypertarget{related-work-and-the-threat-being-addressed}{%
\subsection{2. Related work and the threat being
addressed}\label{related-work-and-the-threat-being-addressed}}

Secure audit logging has long distinguished ordinary local logging from
logs designed to retain protections after the logging machine is
compromised. Schneier and Kelsey's construction explicitly addresses
compromise and uses cryptographic mechanisms beyond an unkeyed local
digest chain. The current PFP store does not implement that construction
or inherit its security properties.
\href{https://www.schneier.com/academic/archives/1999/05/secure_audit_logs_to.html}{Schneier
and Kelsey, 1999}

Certificate Transparency specifies signed tree heads and Merkle
inclusion and consistency proofs for a public logging setting. Those
mechanisms provide a useful reference for identifying what a local
count-and-head record omits: signed producer identity, public witnesses,
and a protocol for comparing log views. The PFP store is not a
Certificate Transparency log.
\href{https://www.rfc-editor.org/rfc/rfc9162.html}{Laurie, Messeri, and
Stradling, 2021}

Reflexion uses remembered linguistic feedback to guide later attempts.
PFP's retained failure records can be inputs to a future feedback
experiment, but record integrity and improved reasoning are different
experimental questions. The present paper studies the former contract
and proposes constraints on reuse.
\href{https://arxiv.org/abs/2303.11366v4}{Shinn et al., 2023}

The immediate threats are accidental mutation, partial truncation, stale
records, and inconsistent reconstruction. A stronger adversary may
rewrite both the entries and the local header. A still stronger
adversary controls the application and any separately retained anchors.
Each level needs an explicit interpretation; a checksum alone does not
resolve them all.

\hypertarget{what-the-store-records}{%
\subsection{3. What the store records}\label{what-the-store-records}}

\texttt{\nolinkurl{record_refusal}} accepts a turn envelope, a
refusal reason, optional prompt parameters, and an optional cochain
projection. It normalizes supported values, assigns an index and
timestamp, and hashes the payload together with the previous entry
digest. Fractions are serialized into rational strings in this storage
path. The following fields are central:

\begin{longtable}[]{@{}
  >{\raggedright\arraybackslash}p{(\columnwidth - 2\tabcolsep) * \real{0.5000}}
  >{\raggedright\arraybackslash}p{(\columnwidth - 2\tabcolsep) * \real{0.5000}}@{}}
\toprule\noalign{}
\begin{minipage}[b]{\linewidth}\raggedright
Field
\end{minipage} & \begin{minipage}[b]{\linewidth}\raggedright
Meaning and limitation
\end{minipage} \\
\midrule\noalign{}
\endhead
\bottomrule\noalign{}
\endlastfoot
\texttt{\nolinkurl{entry_index}} & Nonnegative position in the
ordered chain \\
\texttt{\nolinkurl{turn_envelope}} & Supplied proposal context;
not independently authenticated by storage \\
\texttt{\nolinkurl{refusal_reason}} & Nonempty text identifying
why the caller recorded the event \\
\texttt{\nolinkurl{prompt_parameters}} & Retained parameters
supplied by the caller \\
\texttt{\nolinkurl{cochain_projection}} & Supplied structured
values; not a native calculation performed by the sink \\
\texttt{\nolinkurl{prev_digest}} & Commitment to the preceding
entry, or the genesis marker \\
\texttt{\nolinkurl{entry_digest}} & SHA-256 of the normalized
entry payload \\
Fixed marker field & Must hold the string
\passthrough{\lstinline!None!}; any other value is rejected when an
entry is created or reloaded \\
\end{longtable}

Each payload also carries a record type and the assigned timestamp. A
header records the entry count and current head digest. The exported
\passthrough{\lstinline!TailReceipt!} contains those two values, a
schema identifier, the same fixed marker, and a SHA-256 digest over
those four values. \texttt{\nolinkurl{emit_tail_receipt}} reads
the sink's cached entries and head; it does not reread the file or
recompute the entries. A caller comparing persisted snapshots must first
reopen the store or run \texttt{\nolinkurl{verify_integrity}}
successfully, then emit the observed receipt. This sequence assumes one
writer and controlled custody during the check. The sink can check
consistency of its entries; it cannot determine whether the caller
honestly described an external event.

This distinction matters for mathematical values. A caller can store
\texttt{\{"phi1": "0"\}} as text. Its presence and hash
do not make it a native functional measurement. A record intended to
support such a claim must also identify the actual verifier, input
bytes, output bytes, and source/binary bindings, and a separate
evaluator must check that relationship.

The store does not enforce a general data-handling policy. The caller
must decide which prompts, responses, and parameters it is authorized to
retain before calling it. Rejection is not a retention license. The
sink's fixed \texttt{\nolinkurl{authority_effect}} marker is the
literal string \passthrough{\lstinline!"None"!}; it is neither JSON
\passthrough{\lstinline!null!} nor producer authentication. Sensitive
fields should be excluded, transformed, or handled by an appropriate
authorized storage mechanism. A redacted record can support only the
evidence claim that its remaining information permits.

\hypertarget{logical-append-and-physical-persistence}{%
\subsection{4. Logical append and physical
persistence}\label{logical-append-and-physical-persistence}}

The public API adds entries in increasing order. Physical persistence is
different: \texttt{\nolinkurl{_write_file}} serializes the header
and all retained entries to a temporary file, then replaces the
destination with \passthrough{\lstinline!os.replace!}. Consequently
``append-only'' describes the API's logical history. The file is
rewritten on every append.

This design is simple to inspect and supports restart verification, but
has clear operational limits. It retains the entries in memory, writes
the accumulated contents again, and provides no concurrent-writer lock.
Two independent writers can operate from stale snapshots and interfere
through the same temporary path. There is no explicit
\passthrough{\lstinline!fsync!} call in this path. Process restart tests
therefore do not demonstrate persistence through every power failure or
storage-device fault.

There is also a failed-write boundary before a restart:
\texttt{\nolinkurl{record_refusal}} advances the in-memory entry
list and head before \texttt{\nolinkurl{_write_file}} finishes.
If replacement raises an exception, the caller receives no successful
append, but the old sink object can emit a receipt for an entry absent
from the persisted file. The new package regression injects that one
failure and observes this divergence. Discard or successfully reload
that object and reconcile the attempted operation before retrying; an
exception is not a rollback guarantee. This is a fault-injection
example, not a crash campaign.

If each record has a bounded size B and N records are added
individually, the accumulated entry bytes written are of order B times
N(N+1)/2, plus headers, while a single final snapshot occupies order BN
bytes. This simple count assumes no batching or compaction and is not a
measured benchmark. Bound memory, record size, total retained bytes, and
write failures separately.

The encoding helper uses ordinary compact
\passthrough{\lstinline!json.dumps!} with sorted keys. Its docstring
calls that RFC 8785, but the implementation does not establish the full
serialization and Unicode rules required by that specification. This
paper calls it the sink's sorted-JSON encoding and binds its source.
Interoperability claims require tests against the actual profile and
supported value domain. For example, Python's code-point key order
differs from JCS's UTF-16 code-unit order for some supplementary
characters, and the baseline encoder leaves Python's permissive
non-finite-number option enabled. JCS requires rejection of NaN and
infinity. These are concrete differences, not an assumption that sorted
keys provide standards conformance.
\href{https://www.rfc-editor.org/rfc/rfc8785.html}{RFC 8785}

Normalization is also lossy: \texttt{Fraction(1, 3)}
and the string \texttt{\nolinkurl{"1/3"}} become the same stored
value; integer key \passthrough{\lstinline!1!} and string key
\passthrough{\lstinline!"1"!} collide after key stringification. The
sink therefore retains normalized proposal information, not an injective
encoding of every caller value. Preserve separately authorized raw bytes
and their type/encoding contract when exact original reconstruction is
required.

The returned record is a frozen dataclass, but nested mappings are not
thereby made transitively immutable. Consumers should reconstruct
evidence from verified persisted bytes and avoid treating an in-memory
tuple of records as an immutable storage medium. These limitations do
not invalidate a scoped restart check; they define its scope.

\hypertarget{a-retained-anchor-integrity-argument}{%
\subsection{5. A retained-anchor integrity
argument}\label{a-retained-anchor-integrity-argument}}

Let P\_i be the normalized payload of entry i, including its predecessor
digest, and let h\_i be SHA-256 of that payload's encoded bytes. The
header claims a length n and a head h\_(n-1). Verification reparses each
entry, checks its index and predecessor link, recomputes h\_i, and
compares the resulting count and head to the header. The exercised
empty-file, chain, count, and head checks produce typed integrity or
truncation errors. This is not a claim that every malformed Python or
JSON value follows that error taxonomy.

Under the assumptions that the encoding is deterministic for admitted
values, the parser checks the relevant fields, and a changed payload
cannot be made to collide with the expected digest, changing a protected
entry payload while keeping its original downstream anchor is
detectable. Reordering or removing an interior record also breaks an
index, predecessor link, or the expected head. This is a conditional
integrity argument, not a formal verification of all parser and
filesystem behavior.

The location of the anchor matters. If an attacker removes the last
record but leaves the header unchanged, the local count/head comparison
detects truncation. If the attacker rewrites the header to match a
shorter consistent prefix, that prefix can pass local verification. If
the attacker rewrites every entry and recalculates every digest, local
self-consistency can also hold.

Retaining a prior \passthrough{\lstinline!TailReceipt!} independently
enables a caller to ask whether a reconstructed file has the same
ordered normalized-entry count and head. It does not establish raw-file
or complete header identity. For example, the sink normalizes a refusal
reason through \passthrough{\lstinline!str!}, so the JSON number
\passthrough{\lstinline!42!} and string \passthrough{\lstinline!"42"!}
can produce the same protected payload; header genesis metadata and
unknown header fields are not committed by this comparison. The baseline
loader also does not validate the header's
\passthrough{\lstinline!schema!} value. Replacing that value can leave
the count/head comparison unchanged; an emitted receipt's schema does
not retrospectively validate the file header. A caller requiring exact
file bytes must retain and check an independent raw-file digest as well.
The current sink exposes the receipt, but it does not automatically
retrieve an independently trusted copy or perform a signed freshness
protocol. The consumer must preserve the anchor and compare count and
head. A valid later extension intentionally has a different head;
proving ancestry of that extension requires inspecting the chain against
the retained prefix. Unequal heads alone show only that the two
snapshots differ.

\clearpage
\begin{longtable}[]{@{}
  >{\raggedright\arraybackslash}p{(\columnwidth - 4\tabcolsep) * \real{0.3500}}
  >{\raggedright\arraybackslash}p{(\columnwidth - 4\tabcolsep) * \real{0.1700}}
  >{\raggedright\arraybackslash}p{(\columnwidth - 4\tabcolsep) * \real{0.4800}}@{}}
\toprule\noalign{}
\begin{minipage}[b]{\linewidth}\raggedright
Persisted modification
\end{minipage} & \begin{minipage}[b]{\linewidth}\raggedright
Baseline reload
\end{minipage} & \begin{minipage}[b]{\linewidth}\raggedright
Comparison with retained count/head
\end{minipage} \\
\midrule\noalign{}
\endhead
\bottomrule\noalign{}
\endlastfoot
Change a protected normalized payload, keep old digest & Rejects & No
verified new snapshot \\
Remove suffix, keep old header & Rejects & No verified new snapshot \\
Remove suffix and write a matching header & Can pass & Different from
the retained longer snapshot \\
Add a valid next entry & Passes & Different; inspect the retained prefix
to establish extension \\
Change unbound header metadata or JSON whitespace & Can pass & Can
match; raw-file digest can differ \\
\end{longtable}

The hash assumption concerns replacement of protected encoded payloads
against a retained digest. It does not turn normalizations that already
map different inputs to the same payload into detectable changes, and it
does not authenticate the anchor's custodian. A self-digest of a receipt
can be recomputed by anyone who can replace that receipt.

The same boundary applies to dates. A timestamp records a value written
by the application. The store does not authenticate a clock or prove
that the event happened at that time.

\hypertarget{worked-example-two-refusals-and-three-modifications}{%
\subsection{6. Worked example: two refusals and three
modifications}\label{worked-example-two-refusals-and-three-modifications}}

Use a fresh writable directory and synthetic, nonsensitive data. The
following example calls the actual storage API:

\begin{lstlisting}[language=Python]
from pathlib import Path
from symbol_game.infrastructure.persistence import cemetery_sink

CemeterySink = cemetery_sink.CemeterySink

path = Path("fresh-failure-ledger") / "failures.jsonl"
sink = CemeterySink(path)
sink.record_refusal(
    {"task_id": "toy-1", "candidate": "edit README.md"},
    "PROTECTED_PATH",
    prompt_parameters={"transport": "scripted_fixture"},
)
sink.record_refusal(
    {"task_id": "toy-1", "candidate": "use an absent route"},
    "NO_NATIVE_PORT",
    prompt_parameters={"transport": "scripted_fixture"},
)
anchor = sink.emit_tail_receipt()
reopened = CemeterySink(path)
observed = reopened.emit_tail_receipt()
assert (observed.entry_count, observed.head_digest) == (
    anchor.entry_count, anchor.head_digest
)
\end{lstlisting}

The example compares the normalized-entry snapshot after reopening; it
does not run a code ownership checker or native route. The refusal
reasons are deliberately constructed labels. A real campaign must bind
the host decisions from which those reasons were derived.
Timestamp-dependent digests are obtained from execution, not hardcoded
in the manuscript.

In the first modification, change
\texttt{\nolinkurl{PROTECTED_PATH}} in the first persisted entry
while leaving its digest unchanged. Reopening rejects the changed
payload with \passthrough{\lstinline!LedgerIntegrityError!}. In the
second, delete the final entry while retaining the original header.
Reopening rejects the count/head discrepancy with
\passthrough{\lstinline!LedgerTruncationError!}. The
\href{https://github.com/dbccpage/the_symbol_game/blob/52cb5457f6c485ba548cfd292f6012a941ed22cf/symbol_game/tests/test_cemetery_sink.py}{storage
tests} contain these classes of restart and tamper checks, along with
extra-key, malformed UTF-8, invalid-count, and fixed-marker cases.

In the third modification, construct a complete internally consistent
shorter file. It reopens without error, so local integrity is no longer
enough. A caller comparing its count and head with the independently
held two-entry anchor can observe that the file does not match the
retained snapshot. The anchor must actually be held outside the
rewriting adversary's control; saving it beside the file under the same
untrusted custody does not provide that assumption automatically.

This example is small enough for a new reader to inspect every record.
It also shows why the useful output is a specific comparison, such as
``these entries match this retained snapshot,'' rather than a blanket
claim that the storage is immutable.

\hypertarget{from-a-retained-event-to-a-reusable-detector}{%
\subsection{7. From a retained event to a reusable
detector}\label{from-a-retained-event-to-a-reusable-detector}}

A proposed reuse pipeline starts with a query over retained records,
groups events by a declared key, and selects candidates for
reproduction. Grouping by a digest is a conditional identity test for
the particular covered bytes and encoding, under the stated hash
assumption. A hash of a label or normalized entry does not identify the
entire original proposal or its semantic defect. Grouping by a
normalized feature or presentation key establishes only that chosen
equivalence relation. Semantic equivalence needs a separate checker.

The next step reproduces the failure using the original input and
checker bindings. A failure caused by missing infrastructure must remain
a missing infrastructure case. It should not become a detector that
declares similar mathematical witnesses invalid. If the original bytes
were not retained under policy, record the replay limitation rather than
filling them in from memory.

A reproduced defect can become a proposed rule. The rule is tested on
the failing case and on valid near misses before any enforcement action
is considered. Its provenance includes the source records, predicate,
fixture bank, permitted action, and scope. Reuse is a new evaluation
event, with a new receipt. The current sources supply individual checks
for this process; this paper does not claim an integrated automatic
miner with measured discovery yield.

The
\href{https://github.com/dbccpage/the_symbol_game/blob/52cb5457f6c485ba548cfd292f6012a941ed22cf/symbol_game/domain/immune/clonal_selection.py}{clonal-selection
module} illustrates one such check. It compiles a candidate regular
expression and compares it with benign self-panel examples. A matching
clone produces a kill-condition record naming both clone and parent. An
invalid expression returns
\texttt{\nolinkurl{INVALID_MATCH_PATTERN}}; a Boolean ``not
killed'' value must not be misread as admission when that refusal is
present. A clean result means only no match on the supplied panel. The
method returns a decision and record; it does not mutate the frozen
parent or clone, remove them from a repository, or demonstrate that a
host applied the proposed purge. That consumer action is a separate
integration obligation. Pattern execution also needs appropriate
input-size and resource limits; a syntactically valid regular expression
is not a bounded runtime guarantee.

The
\href{https://github.com/dbccpage/the_symbol_game/blob/52cb5457f6c485ba548cfd292f6012a941ed22cf/symbol_game/tests/test_immune_defect_regressions.py}{self-panel
regression} identifies 40 distinct base contents inside 1,000 stored
examples. A publication must report that effective diversity, rather
than presenting synthetic repetitions as independent coverage. Adding
new domains needs new valid examples and a renewed check.

\hypertarget{the-retained-detector-candidate-catalog}{%
\subsubsection{7.1 The retained detector-candidate
catalog}\label{the-retained-detector-candidate-catalog}}

The repository retains a JSON file of harvested detector candidates,
\href{https://github.com/dbccpage/the_symbol_game/blob/52cb5457f6c485ba548cfd292f6012a941ed22cf/_intelligence/_antibodies/pharmacy_catalog.json}{\texttt{\nolinkurl{_intelligence/_antibodies/pharmacy_catalog.json}}}.
Its inherited name, pharmacy catalog, is a label for a file of records;
the file is not an active filter. At the audited commit it holds 61
entries with 14 distinct \texttt{\nolinkurl{failure_class}}
values. Its \texttt{\nolinkurl{pharmacy_status}} field is
\passthrough{\lstinline!candidate!} for 56 entries and
\passthrough{\lstinline!active!} for 5.

\begin{longtable}[]{@{}
  >{\raggedright\arraybackslash}p{(\columnwidth - 4\tabcolsep) * \real{0.5900}}
  >{\raggedright\arraybackslash}p{(\columnwidth - 4\tabcolsep) * \real{0.2300}}
  >{\raggedright\arraybackslash}p{(\columnwidth - 4\tabcolsep) * \real{0.1800}}@{}}
\toprule\noalign{}
\begin{minipage}[b]{\linewidth}\raggedright
\texttt{\nolinkurl{failure_class}} value
\end{minipage} & \begin{minipage}[b]{\linewidth}\raggedright
Entries
\end{minipage} & \begin{minipage}[b]{\linewidth}\raggedright
\texttt{\nolinkurl{pharmacy_status}}
\end{minipage} \\
\midrule\noalign{}
\endhead
\bottomrule\noalign{}
\endlastfoot
\passthrough{\lstinline!defect:defect:0!} to
\passthrough{\lstinline!defect:defect:4!} (five classes) & 9 each, 45 in
total & candidate \\
\texttt{\nolinkurl{defect:h_boundary}} & 4 & candidate \\
\texttt{\nolinkurl{general_obstruction}} & 3 & candidate \\
\texttt{\nolinkurl{architecture:clean_arch_inward_dependency_violation}}
& 3 & candidate \\
\texttt{\nolinkurl{architecture:circular_dependency}} & 1 &
candidate \\
\texttt{\nolinkurl{algebraic_defect:energy_projector_sign_inversion}}
& 1 & active \\
\texttt{\nolinkurl{algebraic_defect:minimizer_bitmask_convention_inversion}}
& 1 & active \\
\texttt{\nolinkurl{resource_exhaustion:reasoning_trace_exceeds_completion_window}}
& 1 & active \\
\texttt{\nolinkurl{syntax_error:unterminated_string_literal}} &
1 & active \\
\texttt{\nolinkurl{type_error:unsupported_operand_int_and_list}}
& 1 & active \\
\end{longtable}

Each entry has 15 fields, including
\texttt{\nolinkurl{negative_fixture}},
\texttt{\nolinkurl{positive_fixture}}, and
\texttt{\nolinkurl{evidence_sha256}}; every entry has all three.
The fixtures are short strings. In 49 entries they are placeholder
comments: a \texttt{\# Triggering defect:} line that
repeats the defect label, and
\texttt{\# Repaired exact state}. The other 12 entries
hold a descriptive phrase or a code fragment of at most four lines. The
file binds no fixture to a checker, to input bytes, or to a test that
reproduces the failure, and it records no reproduction or near-miss
result.

Each \texttt{\nolinkurl{evidence_sha256}} value is a distinct
64-digit hexadecimal SHA-256 digest. The file does not record which
bytes a digest covers or which code wrote an entry. The two harvesting
implementations at the audited commit hash short identifying strings:
\href{https://github.com/dbccpage/the_symbol_game/blob/52cb5457f6c485ba548cfd292f6012a941ed22cf/moscow/domain/immune/seeder.py}{one}
joins the defect label with the run and task identifiers, and
\href{https://github.com/dbccpage/the_symbol_game/blob/52cb5457f6c485ba548cfd292f6012a941ed22cf/operations/domain/symbol_game_bridge/pharmacy_seeder.py}{the
other} also prefixes the entry identifier. For 19 entries, namely 15 of
the \passthrough{\lstinline!defect:defect:N!} entries and the four
\texttt{\nolinkurl{defect:h_boundary}} entries, the stored value
equals the SHA-256 of the defect label joined with a run identifier and
task identifier that appear in run records retained at that commit. None
of these inputs is a reproduction receipt. As with the head digest in
Section 5, such a value supports comparison only against an
independently retained anchor. It does not make an entry immutable, and
it does not show that a reproduction exists.

No code at the audited commit matches incoming proposals against these
fixtures, executes a fixture, or derives a repair from a positive
fixture. The nearest code found, a
\href{https://github.com/dbccpage/the_symbol_game/blob/52cb5457f6c485ba548cfd292f6012a941ed22cf/external_application/prompt_doctor/use_cases/repair_service.py}{prompt-diagnosis
service}, tests whether an entry's
\texttt{\nolinkurl{target_pattern}} occurs as a substring of a
prompt; it reads neither fixture. The counts above therefore describe
the file. They are not measures of detector coverage or effectiveness.
Admitting any entry as a detector would require the reproduction and
near-miss steps described above.

\hypertarget{monitoring-overactive-rules-and-decay}{%
\subsection{8. Monitoring overactive rules and
decay}\label{monitoring-overactive-rules-and-decay}}

The
\href{https://github.com/dbccpage/the_symbol_game/blob/52cb5457f6c485ba548cfd292f6012a941ed22cf/symbol_game/domain/immune/treg.py}{regulatory
supervisor} computes a toxicity index from trigger records. With the
explicit trigger interface, each supplied record states whether the
rule's pattern matched and whether the second signal of the two-signal
gate was also present. The index is the number of pattern-only records
divided by the number of supplied records. At the default threshold 1/5,
toxicity at or above the threshold synthesizes a record proposing a
downgrade to \texttt{\nolinkurl{DIAGNOSTIC_ONLY}}. Synthesis alone
does not update the rule or host policy. Five records containing one
pattern-only case give exactly 1/5 and therefore meet the threshold. One
such case among six gives 1/6 and remains below it.

These are policy arithmetic, not an estimated population false-positive
rate. The trigger labels must themselves be justified. They describe
signal presence, not an independent oracle judgment that a proposal is
valid or suppressed. Records with neither signal still enter the
explicit-interface denominator. For one pattern-only record, two
dual-signal records, and three no-signal records, this interface reports
1/6; the fraction conditional on a pattern match would instead be 1/3.
Neither denominator is automatically the valid, eligible-proposal
denominator of PFP III's false-suppression experiment. An empty trigger
list produces zero by convention, not evidence of safety. A different
input interface using an unlabelled prompt corpus has different
information available: every match counts as pattern-only, so the score
is 0 when nothing matches and 1 otherwise. That score cannot be
substituted for an oracle-labeled error rate. The regression suite makes
this distinction observable.

Affinity decay provides another bounded mechanism. For idle count t and
positive fixed interval tau and a nonnegative initial score, the
implementation halves a rule's score once for each newly completed
interval: a call that moves the idle count from t to t' applies the
factor (1/2)\^{}k with k = floor(t'/tau) − floor(t/tau). Twelve idle
runs with tau=5 multiply the starting score by 1/4, whether recorded in
one call or batches of 3, 4, and 5. This batching invariance follows
from telescoping the differences of completed interval counts and is
checked by the tests. It assumes tau remains fixed and no stimulation,
reset, or manual score update occurs between batches. Activity expiry is
a separate policy: the rule is marked inactive once its idle count
reaches its configured \texttt{\nolinkurl{half_life_runs}} value.

Neither thresholding nor decay establishes that a rule improves
reasoning. They regulate a proposed detector repertoire. Effectiveness
and safety require the held-out, independently labeled experiment
specified in PFP III. A retained failure can be valuable for diagnosis
even when the proposed detector derived from it is rejected.

\hypertarget{reproducibility-and-evidence-classification}{%
\subsection{9. Reproducibility and evidence
classification}\label{reproducibility-and-evidence-classification}}

Distinguish the manuscript's baseline from a current-checkout regression
run. To reproduce the historical implementation, use the pinned commit
in a separate checkout and record the exact source hashes. Copy the
revised paper package to the corresponding deposit path there if its
added tests are absent at that historical commit. To test the current
branch, report that checkout's own source hashes and any changed
contracts instead of labelling the result a baseline replay.

Collect and run the focused suites from the root of a full repository
checkout, using Python 3.12 (on Windows,
\texttt{py -3.12} in place of
\passthrough{\lstinline!python!}):

\begin{lstlisting}
python -m pytest symbol_game/tests/test_cemetery_sink.py --collect-only -q
python -m pytest symbol_game/tests/test_cemetery_sink.py -q
python -m pytest symbol_game/tests/test_immune_defect_regressions.py --collect-only -q
python -m pytest symbol_game/tests/test_immune_defect_regressions.py -q
\end{lstlisting}

The first suite tests the production sink; the second targets concrete
rule maintenance defects and boundary cases. Use an explicit writable
temporary location where required. Record commands, exits, interpreter,
source hashes, and exact test outcomes. No provider call or model
download is required for these focused checks.

The package test
\texttt{\nolinkurl{code/tests/test_006_worked_example.py}}
executes the Section 6 block as printed, the three modifications
described there, and the number-versus-string reason case from Section 5
against the repository sink. It also exercises unbound header metadata,
whitespace normalization, cached receipt emission, failed replacement,
and mutation of nested mappings. These cases establish the specific
evidence boundaries described above. It needs the full repository
checkout and skips, stating that reason, when the sink source is absent.
It writes ledgers only under pytest's temporary directory. The package
test \texttt{\nolinkurl{code/tests/test_006_catalog_file.py}}
recomputes the Section 7.1 counts, fields, placeholder fixtures, and
digest reconstructions from the catalog bytes read through Git at the
audited commit. The Git path requires that historical object and a
\passthrough{\lstinline!git!} executable; the test falls back to the
identical package-local catalog snapshot when Git or that object is
unavailable. The snapshot is bound to the baseline path, Git blob, and
SHA-256 in \texttt{\nolinkurl{code/fixtures/README.md}}; the tests
check its digest before using it. This makes the catalog arithmetic
rerunnable from the deposit alone, without presenting the catalog as an
executable detector bank.

The paper-local
\texttt{\nolinkurl{code/tests/test_006_epistemic_immune.py}}
defines independent toy classes. It is retained as an educational
illustration and is not evidence that the production sink or immune
pipeline behaves the same way. Its copied positive/zero-defect rule also
differs from the actual two-signal gate. It must not supply production
integration counts.

Current execution receipts must identify which checks actually ran.
Historical receipts and stored outputs remain historical. A successful
restart test establishes the exercised process/reconstruction behavior;
no crash-consistency, power-failure, or concurrent-writer campaign is
claimed. There is also no empirical corpus of mined conjectures
presented as a validated result in this paper.

\hypertarget{limitations-and-next-experiments}{%
\subsection{10. Limitations and next
experiments}\label{limitations-and-next-experiments}}

The storage implementation is suitable for bounded, inspectable local
runs with a single writer and controlled custody. Its trust boundary is
too weak for claims of permanent immutable retention or authenticated
multi-party logging. Source inspection also identifies a naming mismatch
between the sorted-JSON encoder and RFC 8785; interoperability should be
resolved before using a different implementation as a verifier.

The proposed reuse procedure has meaningful unresolved work: retain
anchors under an explicit trust policy, bind actual host decisions into
entries, demonstrate a deterministic reproduction path for each mined
case, and measure rule behavior on independent valid near misses. Any
future integrated miner must report both useful discoveries and harmful
suppressions under a frozen evaluation plan.

The practical result is narrower and useful: a failed attempt can become
an inspectable, source-bound research object. Preserving that object
supports reproduction and review while leaving its mathematical meaning
and future utility open to independent evaluation.

This paper assigns no authorization and changes no jurisdiction. It
carries no operational authority or institutional standing.

\hypertarget{verification-for-this-revision}{%
\subsection{Verification for this
revision}\label{verification-for-this-revision}}

The package verifier (\texttt{\nolinkurl{verify_deposit.py}})
checks the manifest and the two-way file census, and
\texttt{tools/rebuild\_release.py --check} checks the
release archive against both manifests. The fresh
\texttt{\nolinkurl{docs/first_publication_validation_2026_10_02.json}}
receipt and accompanying review identify commands, tool versions, source
hashes, and exact exercised scope. Baseline tests were run against
byte-verified materialized sources, not a full-repository checkout.
Catalog arithmetic uses the bound data snapshot. No current-branch
integration or provider trial is inferred. Historical receipts in this
package keep their original dates and scope and are not current
execution evidence.

\hypertarget{references}{%
\subsection{References}\label{references}}

\begin{itemize}
\tightlist
\item
  Schneier, B., and Kelsey, J. (1999).
  \href{https://www.schneier.com/academic/archives/1999/05/secure_audit_logs_to.html}{Secure
  Audit Logs to Support Computer Forensics}. ACM Transactions on
  Information and System Security 2(2), 159--176. DOI
  10.1145/317087.317089.
\item
  Laurie, B., Messeri, E., and Stradling, R. (2021).
  \href{https://www.rfc-editor.org/rfc/rfc9162.html}{Certificate
  Transparency Version 2.0}. RFC 9162 (Experimental), Internet
  Engineering Task Force, December. DOI 10.17487/RFC9162.
\item
  Shinn, N., Cassano, F., Berman, E., Gopinath, A., Narasimhan, K., and
  Yao, S. (2023). \href{https://arxiv.org/abs/2303.11366v4}{Reflexion:
  Language Agents with Verbal Reinforcement Learning}. arXiv preprint
  arXiv:2303.11366, version 4.
\item
  Rundgren, A., Jordan, B., and Erdtman, S. (2020).
  \href{https://www.rfc-editor.org/rfc/rfc8785.html}{JSON
  Canonicalization Scheme (JCS)}. RFC 8785 (Informational), Internet
  Engineering Task Force, June. DOI 10.17487/RFC8785.
\end{itemize}

Implementation links throughout the paper bind the storage and detector
claims to the stated source baseline. They do not replace independently
retained execution or integrity anchors.

\end{document}
